% Part: first-order-logic
% Chapter: tableaux
% Section: proof-theoretic-notions

\documentclass[../../../include/open-logic-section]{subfiles}

\begin{document}

\iftag{FOL}
      {\olfileid{fol}{tab}{ptn}}
      {\olfileid{pl}{tab}{ptn}}
      
\olsection{Bewijstheoretische begrippen}

\begin{editorial}
Deze paragraaf bundelt de definities van de afleidbaarheidsrelatie
en consistentie voor tableaus.
\end{editorial}

\begin{explain}
Zoals we enkele belangrijke semantische begrippen hebben gedefinieerd
(geldigheid, logisch gevolg en vervulbaarheid), definiëren we nu de
overeenkomstige \emph{bewijstheoretische begrippen}. Deze worden niet
gedefinieerd door na te gaan of !!{sentence}s in !!{structure}s waar
zijn, maar door na te gaan of bepaalde gesloten !!{tableau}s bestaan.
Het was een belangrijke ontdekking dat deze begrippen samenvallen. Dit
is de inhoud van de \emph{correctheids-} en
\emph{volledigheidsstellingen}.
\end{explain}


\begin{defn}[Stellingen]
Een !!{sentence}~$!A$ is een \emph{stelling} als er een gesloten
!!{tableau} bestaat voor~$\sFmla{\False}{!A}$. We schrijven
$\Proves !A$ als $!A$ een stelling is en $\Proves/ !A$ als dat niet
zo is.
\end{defn}

\begin{defn}[!!^{derivability}]
!!^a{sentence} $!A$ is \emph{!!{derivable} uit} een verzameling
!!{sentence}s~$\Gamma$, $\Gamma \Proves !A$, dan en slechts dan als er
een eindige verzameling $\{!B_1, \dots, !B_n\} \subseteq \Gamma$ is
en er een gesloten !!{tableau} bestaat voor de verzameling
\[
\{
  \sFmla{\False}{!A},
  \sFmla{\True}{!B_1}, \dots,
  \sFmla{\True}{!B_n}
  \}.
\]
Als $!A$ niet !!{derivable} is uit $\Gamma$, schrijven we
$\Gamma \Proves/ !A$.
\end{defn}

\begin{defn}[Consistentie]
Een verzameling !!{sentence}s~$\Gamma$ heet \emph{inconsistent} dan en
slechts dan als er een eindige verzameling
$\{!B_1, \dots, !B_n\} \subseteq \Gamma$ is en er een gesloten
!!{tableau} bestaat voor de verzameling
\[
\{
  \sFmla{\True}{!B_1}, \dots,
  \sFmla{\True}{!B_n}  
  \}.
\]
Als $\Gamma$ niet inconsistent is, noemen we haar \emph{consistent}.
\end{defn}

\begin{prop}[Reflexiviteit]
\ollabel{prop:reflexivity}
Als $!A \in \Gamma$, dan $\Gamma \Proves !A$.
\end{prop}

\begin{proof}
Als $!A \in \Gamma$, dan is $\{!A\}$ een eindige deelverzameling
van~$\Gamma$ en is het !!{tableau}
\begin{oltableau}
  [\sFmla{\False}{\formula{A}}, just = \TAss
    [\sFmla{\True}{\formula{A}}, just = \TAss,close]
  ]
\end{oltableau}
gesloten.
\end{proof}
  

\begin{prop}[Monotoniciteit]
\ollabel{prop:monotonicity}
Als $\Gamma \subseteq \Delta$ en $\Gamma \Proves !A$, dan $\Delta
\Proves !A$.
\end{prop}

\begin{proof}
Elke eindige deelverzameling van~$\Gamma$ is ook een eindige
deelverzameling van~$\Delta$.
\end{proof}

\begin{prop}[Transitiviteit]
\ollabel{prop:transitivity}
Als $\Gamma \Proves !A$ en $\{!A\} \cup
\Delta \Proves !B$, dan $\Gamma \cup \Delta \Proves !B$.
\end{prop}

\begin{proof}
Als $\{!A\} \cup \Delta \Proves !B$, dan is er een eindige
deelverzameling $\Delta_0 =
\{!C_1, \dots, !C_n\} \subseteq \Delta$ zodat
\begin{align*}
\{\sFmla{\False}{!B}, & \sFmla{\True}{!A}, \sFmla{\True}{!C_1},
\dots, \sFmla{\True}{!C_n}\}
\intertext{een gesloten !!{tableau} heeft. Als $\Gamma \Proves !A$,
  dan zijn er $!D_1$, \dots, $!D_m \subseteq \Gamma$ zodat}
\{\sFmla{\False}{!A}, & \sFmla{\True}{!D_1},
\dots, \sFmla{\True}{!D_m}\}
\end{align*}
een gesloten !!{tableau} heeft.

Beschouw nu het !!{tableau} met de aannamen
\[
\sFmla{\False}{!B},
\sFmla{\True}{!C_1}, \dots, \sFmla{\True}{!C_n},
\sFmla{\True}{!D_1}, \dots, \sFmla{\True}{!D_m}.
\]
Pas op~$!A$ de regel \Cut{} toe. Hierdoor ontstaan twee takken: de ene
bevat $\sFmla{\True}{!A}$, de andere $\sFmla{\False}{!A}$. Op de ene
tak zijn dus alle elementen van
\[
\{\sFmla{\False}{!B}, \sFmla{\True}{!A},
\sFmla{\True}{!C_1}, \dots, \sFmla{\True}{!C_n}\}
\]
beschikbaar. Omdat er voor deze aannamen een gesloten !!{tableau}
bestaat, kunnen we dat aan deze tak vastmaken; elke tak die door
$\sFmla{\True}{!A}$ loopt, sluit. Op de andere tak zijn alle elementen
van
\[
\{\sFmla{\False}{!A}, \sFmla{\True}{!D_1}, \dots,
\sFmla{\True}{!D_m}\}
\]
beschikbaar. We kunnen dus ook de andere kant aanvullen tot een
gesloten !!{tableau}. Dit bewijst $\Gamma \cup \Delta \Proves !B$.
\end{proof}

Dit betekent in het bijzonder dat als $\Gamma \Proves !A$ en $!A
\Proves !B$, dan $\Gamma \Proves !B$. Hieruit volgt ook dat als $!A_1,
\dots, !A_n \Proves !B$ en $\Gamma \Proves !A_i$ voor elke~$i$, dan
$\Gamma \Proves !B$.

\begin{prop}
\ollabel{prop:incons}
$\Gamma$ is inconsistent dan en slechts dan als $\Gamma \Proves !A$
voor elke !!{sentence}~$!A$.
\end{prop}

\begin{proof}
Opgave.
\end{proof}

\tagprob{FOL}
\begin{prob}
Bewijs \olref[fol][tab][ptn]{prop:incons}
\end{prob}
\tagendprob

\tagprob{notFOL}
\begin{prob}
Bewijs \olref[pl][tab][ptn]{prop:incons}
\end{prob}
\tagendprob

\begin{prop}[Compactheid]
  \ollabel{prop:proves-compact}
  \begin{enumerate}
  \item Als $\Gamma \Proves !A$, dan is er een eindige deelverzameling
    $\Gamma_0 \subseteq \Gamma$ zodat $\Gamma_0 \Proves !A$.
  \item Als elke eindige deelverzameling van~$\Gamma$ consistent is,
    dan is $\Gamma$ consistent.
  \end{enumerate}
\end{prop}

\begin{proof}
  \begin{enumerate}
    \item Als $\Gamma \Proves !A$, dan bestaan er een eindige
      deelverzameling $\Gamma_0 = \{!B_1, \dots, !B_n\}$ en een
      gesloten !!{tableau} voor
      \[
      \{\sFmla{\False}{!A}, \sFmla{\True}{!B_1}, \dots, \sFmla{\True}{!B_n}\}
      \]
      Dit !!{tableau} bewijst ook $\Gamma_0 \Proves !A$.
    \item Als $\Gamma$ inconsistent is, dan bestaat er voor een eindige
      deelverzameling $\Gamma_0 = \{!B_1, \dots, !B_n\}$ een gesloten
      !!{tableau} voor
      \[
      \{\sFmla{\True}{!B_1}, \dots, \sFmla{\True}{!B_n}\}
      \]
      Dit gesloten !!{tableau} bewijst dat $\Gamma_0$ inconsistent is.
  \end{enumerate}
\end{proof}

\end{document}
